\ifdefined\gkver\else\def\gkver{student}\fi
\documentclass[\gkver]{gaokaozhenti}
\begin{document}

\chapter{1998年全国}

\begin{enumerate}
\item
%% number: 1
%% paperName: 1998年全国高考第 1 题 5 分
%% typeId: 1
%% type: 单选题
%% chapter: 力物体的平衡
%% point: 三力平衡
%% method: 无
%% score: 5
%% degree: 850
%% duplicateId: 0
%% body:
三段不可伸长的细绳 OA、OB、OC 能承受的最大拉力相同，它们共同悬挂一重物，如图所示，其中 OB 是水平的，A 端、B 端固定。若逐渐增加 C 端所挂物体的质量，则最先断的绳\xzanswer{A}
\onepicture[w=3.5cm]{\includesvg{figs/01}}
\fourchoices[answer=A]
{必定是 OA}
{必定是 OB}
{必定是 OC}
{可能是 OB，也可能是 OC}

%% answer: A
%% memo:
\memoanswer{
由题意知，OA、OB、OC 所受的拉力关系为 $F_{OA}^{2}=F_{OB}^{2}+F_{OC}^{2}$，所以 OA 所受拉力最大，即 OA 先断，故 A 正确。
}

\item
%% number: 2
%% paperName: 1998年全国高考第 2 题 5 分
%% typeId: 2
%% type: 多选题
%% chapter: 气体性质
%% point: 分子动理论
%% method: 无
%% score: 5
%% degree: 850
%% duplicateId: 0
%% body:
下列说法正确的是\xzanswer{AB}
\fourchoices[answer=AB]
{液体中悬浮微粒的布朗运动是作无规则运动的液体分子撞击微粒而引起的}
{物体的温度越高，其分子的平均动能越大}
{物体里所有分子动能的总和叫做物体的内能}
{只有传热才能改变物体的内能}

%% answer: AB
%% memo:
\memoanswer{
布朗运动是液体分子撞击微粒造成的，温度是分子平均动能的标志，故 AB 正确；内能是指物体里所有分子动能与势能的总和，做功和热传递都能改变物体的内能，故 CD 错误。
}

\item
%% number: 3
%% paperName: 1998年全国高考第 3 题 5 分
%% typeId: 2
%% type: 多选题
%% chapter: 机械波
%% point: 波的图象
%% method: 无
%% score: 5
%% degree: 650
%% duplicateId: 0
%% body:
一简谐横波在 $x$ 轴上传播，在某时刻的波形如图所示，已知此时质点 F 的运动方向向下，则\xzanswer{AB}
\onepicture[w=7cm]{\includesvg{figs/03}}
\fourchoices[answer=AB]
{此波朝 $x$ 轴负方向传播}
{质点 D 此时向下运动}
{质点 B 将比质点 C 先回到平衡位置}
{质点 E 的振幅为零}

%% answer: AB
%% memo:
\memoanswer{
根据 $F$ 的振动方向，利用同侧法可以判断，波谷向左移动，波向左传播，同理可判断 $E$、$D$、$C$ 质点向下运动，故 AB 正确；由于波向左传播，$C$ 处波峰向左平移，可知 $B$ 点向上运动，比 $C$ 晚一些回到平衡位置，故 CD 错误。
}

\item
%% number: 4
%% paperName: 1998年全国高考第 4 题 5 分
%% typeId: 2
%% type: 多选题
%% chapter: 磁场
%% point: 左手定则
%% method: 无
%% score: 5
%% degree: 650
%% duplicateId: 0
%% body:
通电矩形导线框 abcd 与无限长通电直导线 MN 在同一平面内，电流方向如图所示，ab 边与 MN 平行。关于 MN 的磁场对线框的作用，下列叙述正确的是\xzanswer{BC}
\onepicture[w=5cm]{\includesvg{figs/04}}
\fourchoices[answer=BC]
{线框有两条边所受的安培力方向相同}
{线框有两条边所受的安培力大小相同}
{线框所受安培力的合力朝左}
{cd 所受安培力对 ab 边的力矩不为零}

%% answer: BC
%% memo:
\memoanswer{
线框四边处在 MN 中电流产生的垂直纸面向里的非均匀磁场中，且离 MN 越远，磁场越弱，当线框通以顺时针电流时，由左手定则可判知四条边所受安培力方向没有相同的，故 A 错误；由于无限长直导线在 ad 边和 bc 边处的磁场是完全相同的，所以这两边所受安培力大小相同，故 B 正确；由于 ab 边处磁感应强度大于 cd 边处磁感应强度，两边受安培力有 $F_{ab}>F_{cd}$，方向相反，合力向左，ad、bc 两边安培力之和为零，故 C 正确；因为四边所受安培力均在线圈平面内，故 cd 所受安培力的力矩应为零，故 D 错误。
}

\item
%% number: 5
%% paperName: 1998年全国高考第 5 题 5 分
%% typeId: 1
%% type: 单选题
%% chapter: 电磁感应
%% point: 电磁感应中图象
%% method: 无
%% score: 5
%% degree: 450
%% duplicateId: 0
%% body:
如图所示，一宽 $40\Ucm$ 的匀强磁场区域，磁场方向垂直纸面向里。一边长为 $20\Ucm$ 的正方形导线框位于纸面内，以垂直于磁场边界的恒定速度 $v=20\Ucms$ 通过磁场区域，在运动过程中，线框有一边始终与磁场区域的边界平行。取它刚进入磁场的时刻 $t=0$，在下列图线中，正确反映感应电流强度随时间变化规律的是\xzanswer{C}
\onepicture[w=3.5cm]{figs/05a.png}
\fourchoices[answer=C, ispicture=true, h=2.4cm]
{figs/05b.png}
{figs/05c.png}
{figs/05d.png}
{figs/05e.png}

%% answer: C
%% memo:
\memoanswer{
线框进、出磁场，线框中磁通量先增后减，磁场方向不变，所以线框中感应电流方向相反；当线框全部进入磁场后，回路中磁通量不变，无感应电流，故 C 正确。
}

\item
%% number: 6
%% paperName: 1998年全国高考第 6 题 5 分
%% typeId: 2
%% type: 多选题
%% chapter: 原子和原子核
%% point: 天然放射性
%% method: 无
%% score: 5
%% degree: 650
%% duplicateId: 0
%% body:
天然放射性元素 \ce{^{232}_{90}Th}（钍）经过一系形 $\alpha$ 衰变和 $\beta$ 衰变之后，变成 \ce{^{208}_{82}Pb}（铅）。下列论断中正确的是\xzanswer{BD}
\fourchoices[answer=BD]
{铅核比钍核少 24 个中子}
{铅核比钍核少 8 个质子}
{衰变过程中共有 4 次 $\alpha$ 衰变和 8 次 $\beta$ 衰变}
{衰变过程中共有 6 次 $\alpha$ 衰变和 4 次 $\beta$ 衰变}

%% answer: BD
%% memo:
\memoanswer{
铅核的中子数 $208-82=126$，钍核的中子数 $232-90=142$，铅核比钍核少 $142-126=16$ 个中子，故 A 错误；铅核比钍核少 $90-82=8$ 个质子，故 B 正确；衰变的核反应方程为 $\ce{^{232}_{90}Th -> ^{208}_{82}Pb + x^{4}_{2}He + y^{0}_{-1}e}$，由质量数守恒得 $232=208+4x$，由电荷数守恒得 $90=82+2x-y$，解得 $x=6$，$y=4$，故 D 正确。
}

\item
%% number: 7
%% paperName: 1998年全国高考第 7 题 5 分
%% typeId: 1
%% type: 单选题
%% chapter: 电场
%% point: 静电感应
%% method: 无
%% score: 5
%% degree: 650
%% duplicateId: 0
%% body:
一金属球，原来不带电，现沿球的直径的延长线放置一均匀带电的细杆 MN，如图所示，金属球上感应电荷产生的电场在球内直径上 a、b、c 三点的场强大小分别为 $E_{a}$、$E_{b}$、$E_{c}$，三者相比\xzanswer{C}
\onepicture[w=6cm]{figs/07.png}
\fourchoices[answer=C]
{$E_{a}$ 最大}
{$E_{b}$ 最大}
{$E_{c}$ 最大}
{$E_{a}=E_{b}=E_{c}$}

%% answer: C
%% memo:
\memoanswer{
由静电平衡知，a、b、c 三点的合场强为零，因此三点由带电细杆 MN 产生的场强与金属球上感应电荷产生的场强大小相等、方向相反，细杆 MN 产生的电场在 c 点的场强最大，则金属球上感应电荷产生的电场在 c 点的场强 $E_{c}$ 最大，故 C 正确。
}

\item
%% number: 8
%% paperName: 1998年全国高考第 8 题 5 分
%% typeId: 1
%% type: 单选题
%% chapter: 光学
%% point: 光的折射
%% method: 无
%% score: 5
%% degree: 650
%% duplicateId: 0
%% body:
一细束平行光经玻璃三棱镜折射后分解为互相分离的三束光，分别照射到相同的金属板 a、b、c 上，如图所示。已知金属板 b 有光电子放出，则可知\xzanswer{D}
\onepicture[w=7cm]{figs/08.png}
\fourchoices[answer=D]
{板 a 一定不放出光电子}
{板 a 一定放出光电子}
{板 c 一定不放出光电子}
{板 c 一定放出光电子}

%% answer: D
%% memo:
\memoanswer{
同一介质对不同频率的色光折射率不同，频率越高折射率越大，通过三棱镜后偏折得越厉害，由题图可知，分别射向金属板 a、b、c 的三束光，频率越来越高，a、b、c 是相同的金属板，发生光电效应极限频率相同，b 板有光电子放出，则 c 板一定有光电子放出，a 板不一定有光电子放出，故 D 正确。
}

\item
%% number: 9
%% paperName: 1998年全国高考第 9 题 5 分
%% typeId: 1
%% type: 单选题
%% chapter: 原子和原子核
%% point: 玻尔模型
%% method: 无
%% score: 5
%% degree: 650
%% duplicateId: 0
%% body:
处于基态的氢原子在某单色光束照射下，只能发出频率为 $\nu_{1}$、$\nu_{2}$、$\nu_{3}$ 的三种光，且 $\nu_{1}<\nu_{2}<\nu_{3}$，则该照射光的光子能量为\xzanswer{C}
\fourchoices[answer=C]
{$h\nu_{1}$}
{$h\nu_{2}$}
{$h\nu_{3}$}
{$h\left(\nu_{1}+\nu_{2}+\nu_{3}\right)$}

%% answer: C
%% memo:
\memoanswer{
处于基态的原子要发光，必须先吸收一定的能量 $E$，如果是用光照射来提供这个能量，则 $E=h\nu$，使之处于激发态，由于激发态能量高，原子不稳定，就会向低能级跃迁，从而发出一系列频率的光子，但这些光子的频率绝不会大于 $\nu$，且必有一种频率等于 $\nu$，由题意知，该氢原子受激后只能发出频率为 $\nu_{1}$、$\nu_{2}$、$\nu_{3}$ 的三种光，且 $\nu_{1}<\nu_{2}<\nu_{3}$，即最高频率是 $\nu_{3}$，那么照射光频率必是 $\nu_{3}$，光子能量是 $h\nu_{3}$，故 C 正确。
}

\item
%% number: 10
%% paperName: 1998年全国高考第 10 题 5 分
%% typeId: 2
%% type: 多选题
%% chapter: 运动和力
%% point: 动量和能量
%% method: 无
%% score: 5
%% degree: 550
%% duplicateId: 0
%% body:
在光滑水平面上，动能为 $E_{0}$、动量的大小为 $p_{0}$ 的小钢球 1 与静止小钢球 2 发生碰撞，碰撞前后球 1 的运动方向相反。将碰撞后球 1 的动能和动量的大小分别记为 $E_{1}$、$p_{1}$，球 2 的动能和动量的大小分别记为 $E_{2}$、$p_{2}$，则必有\xzanswer{ABD}
\fourchoices[answer=ABD]
{$E_{1}<E_{0}$}
{$p_{1}<p_{0}$}
{$E_{2}>E_{0}$}
{$p_{2}>p_{0}$}

%% answer: ABD
%% memo:
\memoanswer{
两个钢球组成的系统在碰撞过程中动量守恒，设钢球 1 初动量的方向为正方向，由动量守恒定律得 $p_{0}=-p_{1}+p_{2}$，碰后球 2 的动量 $p_{2}=p_{0}+p_{1}$，可见 $p_{2}>p_{0}$，故 D 正确；单从动量方面分析，$p_{1}$ 可以大于 $p_{0}$，若如此必有碰后系统的机械能增加，但对于这一具体问题来说碰撞过程没有其他形式的能向机械能转化，机械能只能向其他形式能转化，因此 $E_{1}+E_{2}\leq E_{0}$，必有 $E_{1}<E_{0}$，$E_{2}<E_{0}$，故 A 正确 C 错误；由 $p=\sqrt{2mE_{k}}$ 结合 $E_{1}<E_{0}$ 得 $p_{1}<p_{0}$，故 B 正确。
}

\item
%% number: 11
%% paperName: 1998年全国高考第 11 题 5 分
%% typeId: 2
%% type: 多选题
%% chapter: 电路
%% point: 串并联电路
%% method: 无
%% score: 5
%% degree: 450
%% duplicateId: 0
%% body:
图示的两种电路中，电源相同，各电阻器阻值相等，各电流表的内阻相等且不可忽略不计。电流表 A$_{1}$、A$_{2}$、A$_{3}$ 和 A$_{4}$ 读出的电流值分别为 $I_{1}$、$I_{2}$、$I_{3}$ 和 $I_{4}$。下列关系式中正确的是\xzanswer{BD}
\onepicture[w=8cm]{figs/11.png}
\fourchoices[answer=BD]
{$I_{1}=I_{3}$}
{$I_{1}<I_{4}$}
{$I_{2}=2I_{1}$}
{$I_{2}<I_{3}+I_{4}$}

%% answer: BD
%% memo:
\memoanswer{
因为上图中两支路电阻不等（$I_{1}$ 支路多了电流表内阻），因此分流不等，有 $I_{2}>2I_{1}$，故 C 错误；设各电流表的内阻均为 $r$，上图中回路的总电阻 $R_{1}=\frac{(R+r)\cdot R}{2R+r}+r$，下图中回路的总电阻 $R_{2}=\frac{R+r}{2}$，$R_{1}-R_{2}=\frac{r(3R+r)}{2(2R+r)}>0$，即 $R_{1}>R_{2}$，所以 $I_{2}<I_{3}+I_{4}$，故 D 正确；因 $I_{1}<\frac{I_{2}}{2}$，$I_{3}=\frac{I_{3}+I_{4}}{2}$，$I_{2}<I_{3}+I_{4}$，所以 $I_{1}<I_{3}$，故 A 错误；因 $I_{1}<\frac{I_{2}}{2}$，$I_{4}=\frac{I_{3}+I_{4}}{2}$，$I_{2}<I_{3}+I_{4}$，所以 $I_{1}<I_{4}$，故 B 正确。
}

\item
%% number: 12
%% paperName: 1998年全国高考第 12 题 5 分
%% typeId: 2
%% type: 多选题
%% chapter: 机械振动
%% point: 单摆
%% method: 无
%% score: 5
%% degree: 450
%% duplicateId: 0
%% body:
图中两单摆摆长相同，平衡时两摆球刚好接触。现将摆球 A 在两摆线所在平面内向左拉开一小角度后释放，碰撞后，两摆球分开各自做简谐运动，以 $m_{A}$、$m_{B}$ 分别表示摆球 A、B 的质量，则\xzanswer{CD}
\onepicture[w=3.5cm]{\includesvg{figs/12}}
\fourchoices[answer=CD]
{如果 $m_{A}>m_{B}$，下一次碰撞将发生在平衡位置右侧}
{如果 $m_{A}<m_{B}$，下一次碰撞将发生在平衡位置左侧}
{无论两摆球的质量之比是多少，下一次碰撞都不可能在平衡位置右侧}
{无论两摆球的质量之比是多少，下一次碰撞都不可能在平衡位置左侧}

%% answer: CD
%% memo:
\memoanswer{
单摆做简谐运动的周期为 $T=2\pi\sqrt{\frac{L}{g}}$，与摆球的质量无关，因此两单摆周期相同，碰后经过 $\frac{1}{2}T$ 都将回到最低点再发生碰撞，下一次碰撞一定发生在平衡位置，不可能在平衡位置左侧或右侧，故 CD 正确。
}

\item
%% number: 13
%% paperName: 1998年全国高考第 13 题 4 分
%% typeId: 3
%% type: 填空题
%% chapter: 交流电
%% point: 变压器
%% method: 无
%% score: 4
%% degree: 650
%% duplicateId: 0
%% body:
一理想变压器，原线圈匝数 $n_{1}=1100$，接在电压 $220\UV$ 的交流电源上。当它对 11 只并联的“$36\UV$，$60\UW$”灯泡供电时，灯泡正常发光。由此可知该变压器副线圈的匝数 $n_{2}=$\tkanswer[2cm]{180}，通过原线圈的电流 $I_{1}=$\tkanswer[2cm]{3}$\UA$。

%% answer: 180，3
\jdanswer{
180，3
}

%% memo:
\memoanswer{
因 11 只并联的“36 V，60 W”灯泡正常发光，所以副线圈输出的电压必为 $36\UV$，由 $\frac{U_{1}}{U_{2}}=\frac{n_{1}}{n_{2}}$ 得 $n_{2}=\frac{U_{2}}{U_{1}}n_{1}=\frac{36\times1100}{220}=180$；每只灯泡正常发光时的电流由 $P_{2}'=I_{2}'U_{2}'$ 得 $I_{2}'=\frac{P_{2}'}{U_{2}'}=\frac{60}{36}\UA=\frac{5}{3}\UA$，副线圈中的电流 $I_{2}=11I_{2}'=\frac{55}{3}\UA$，由 $\frac{I_{1}}{I_{2}}=\frac{n_{2}}{n_{1}}$ 得 $I_{1}=\frac{n_{2}}{n_{1}}I_{2}=\frac{180}{1100}\times\frac{55}{3}\UA=3\UA$。
}

\item
%% number: 14
%% paperName: 1998年全国高考第 14 题 4 分
%% typeId: 3
%% type: 填空题
%% chapter: 光学
%% point: 光的折射
%% method: 无
%% score: 4
%% degree: 650
%% duplicateId: 0
%% body:
如图所示，一储油圆桶，底面直径与桶高均为 $d$。当桶内无油时，从某点 A 恰能看到桶底边缘上的某点 B。当桶内油的深度等于桶高一半时，在 A 点沿 AB 方向看去，看到桶底上的 C 点，C、B 相距 $\frac{1}{4}d$。由此可得油的折射率 $n=$\tkanswer[2.5cm]{$\frac{\sqrt{10}}{2}$}；光在油中传播的速度 $v=$\tkanswer[2.8cm]{$6\sqrt{10}\times10^{7}$}$\Ums$。（结果可用根式表示）
\onepicture[align=right, w=5cm]{\includesvg{figs/14}}

%% answer: $\frac{\sqrt{10}}{2}$，$6\sqrt{10}\times10^{7}$
\jdanswer{
$\frac{\sqrt{10}}{2}$，$6\sqrt{10}\times10^{7}$
}

%% memo:
\memoanswer{
光路图如图所示，由折射定律得 $n=\frac{\sin i}{\sin r}$，又 $\sin i=\sin\angle BOD=\frac{\overline{BD}}{\overline{OB}}=\frac{1}{\sqrt{2}}$，$\sin r=\sin\angle COD=\frac{\overline{CD}}{\overline{OC}}=\frac{\frac{1}{4}}{\sqrt{\left(\frac{1}{2}\right)^{2}+\left(\frac{1}{4}\right)^{2}}}$，解得 $n=\frac{\sqrt{10}}{2}$；光在油中传播的速度为 $v=\frac{c}{n}=\frac{3\times10^{8}}{\frac{\sqrt{10}}{2}}\Ums=6\sqrt{10}\times10^{7}\Ums$。

\onepicture[w=5cm]{figs/14_an.png}
}

\item
%% number: 15
%% paperName: 1998年全国高考第 15 题 4 分
%% typeId: 3
%% type: 填空题
%% chapter: 运动和力
%% point: 牛顿第二定律
%% method: 整体与隔离
%% score: 4
%% degree: 650
%% duplicateId: 0
%% body:
如图，质量为 $2m$ 的物块 A 与水平地面的摩擦可忽略不计，质量为 $m$ 的物块 B 与地面的摩擦系数为 $\mu$。在已知水平推力 $F$ 的作用下，A、B 作加速运动。A 对 B 的作用力为\tkanswer[3.5cm]{$\frac{F+2\mu mg}{3}$}。
\onepicture[align=right, w=7cm]{\includesvg{figs/15}}

%% answer: $\frac{F+2\mu mg}{3}$
\jdanswer{
$\frac{F+2\mu mg}{3}$
}

%% memo:
\memoanswer{
因 A、B 一起加速运动，可看做整体，由牛顿第二定律得 $F-\mu mg=3ma$，解得 $a=\frac{F-\mu mg}{3m}$；隔离 B，水平方向上受摩擦力的作用，有 $F_{f}=\mu mg$，A 对 B 的作用力 $T$，由牛顿第二定律得 $T-\mu mg=ma$，解得 $T=\mu mg+\frac{F-\mu mg}{3}=\frac{F+2\mu mg}{3}$。
}

\item
%% number: 16
%% paperName: 1998年全国高考第 16 题 4 分
%% typeId: 3
%% type: 填空题
%% chapter: 电场
%% point: 库仑定律
%% method: 无
%% score: 4
%% degree: 450
%% duplicateId: 0
%% body:
A、B 两带电小球，A 固定不动，B 的质量为 $m$。在库仑力作用下，B 由静止开始运动。已知初始时，A、B 间的距离为 $d$，B 的加速度为 $a$。经过一段时间后，B 的加速度变为 $\frac{1}{4}a$，此时 A、B 间的距离应为\tkanswer[2cm]{$2d$}。已知此时 B 的速度为 $v$，则在此过程中电势能的减少量为\tkanswer[2.5cm]{$\frac{1}{2}mv^{2}$}。

%% answer: $2d$，$\frac{1}{2}mv^{2}$
\jdanswer{
$2d$，$\frac{1}{2}mv^{2}$
}

%% memo:
\memoanswer{
当 B 的加速度为 $a$ 时，有 $a=k\frac{Q_{A}Q_{B}}{md^{2}}$，其中 $Q_{A}$、$Q_{B}$ 为 A、B 两球所带电量，当 B 的加速度变为 $\frac{a}{4}$ 时，设此时 A、B 两球距离为 $x$，有 $\frac{a}{4}=k\frac{Q_{A}Q_{B}}{mx^{2}}$，解得 $x=2d$；B 球在运动过程中，电场力做的功等于其电势能的减少量，由动能定理知，因只受电场力作用，所以电场力做的功等于其动能的增加量，因此 B 球在此过程中电势能的减少量与其动能的增加量相等，即为 $\frac{1}{2}mv^{2}$。
}

\item
%% number: 17
%% paperName: 1998年全国高考第 17 题 4 分
%% typeId: 3
%% type: 填空题
%% chapter: 电场
%% point: 带电粒子的直线运动
%% method: 无
%% score: 4
%% degree: 350
%% duplicateId: 0
%% body:
来自质子源的质子（初速度为零），经一加速电压为 $800\UkV$ 的直线加速器加速，形成电流强度为 $1\UmA$ 的细柱形质子流。已知质子电荷 $e=1.60\times10^{-19}\UC$。这束质子流每秒打到靶上的质子数为\tkanswer[3cm]{$6.25\times10^{15}$}。假定分布在质子源到靶之间的加速电场是均匀的，在质子束中与质子源相距 $l$ 和 $4l$ 的两处，各取一段极短的相等长度的质子流，其中的质子数分别为 $n_{1}$ 和 $n_{2}$，则 $\frac{n_{1}}{n_{2}}=$\tkanswer[2cm]{2}。

%% answer: $6.25\times10^{15}$，2
\jdanswer{
$6.25\times10^{15}$，2
}

%% memo:
\memoanswer{
电流是在真空中的电子的流动，由于电子不断地被加速，即电子运动越来越快，而电流恒定，即单位时间内通过的电荷量不变，亦即相同时间内通过的电子个数相同，所以相同长度内电子个数不同，因为速度越大，时间越短，则电子个数越少。

由电流强度的定义知 $I=\frac{q}{t}=\frac{eN}{t}$，解得 $N=\frac{It}{e}=\frac{1\times10^{3}\times1}{1.6\times10^{-19}}=6.25\times10^{15}$ 个。

\onepicture[w=9cm]{figs/17_an.png}

柱形质子流示意图如图所示，质子在加速电场中做匀加速直线运动，设在与质子源相距 $l$ 和 $4l$ 两处的速度为 $v_{1}$、$v_{2}$，在极短的长度内，质子的运动可看成匀速运动，质子运动的时间为 $\Delta t$，有 $\Delta l=v\Delta t$，得 $\Delta t=\frac{\Delta l}{v}$，则有 $\Delta l\propto\frac{1}{v}$。

质子加速度的距离为 $x$ 时，由动能定理得 $\frac{1}{2}mv^{2}=eU_{x}=eU\frac{x}{L}$，得 $v=\sqrt{\frac{2eU}{mL}}\cdot\sqrt{x}$。

电流恒定，即单位时间内通过的电荷量不变，通过的电子个数为 $n=\frac{q}{e}=\frac{I\Delta t}{e}=\frac{I\Delta l}{e}\cdot\sqrt{\frac{mL}{2eU}}\cdot\frac{1}{\sqrt{x}}$，故 $\frac{n_{1}}{n_{2}}=\frac{\sqrt{x_{2}}}{\sqrt{x_{1}}}=\frac{\sqrt{4l}}{\sqrt{l}}=\frac{2}{1}=2$。
}

\item
%% number: 18
%% paperName: 1998年全国高考第 18 题 6 分
%% typeId: 4
%% type: 实验题
%% chapter: 交流电
%% point: LC振荡回路
%% method: 无
%% score: 6
%% degree: 650
%% duplicateId: 0
%% body:
在 LC 振荡电路中，如已知电容 $C$，并测得电路的固有振荡周期 $T$，即可求得电感 $L$。为了提高测量精度，需多次改变 $C$ 值并侧得相应的 $T$ 值。现将测得的六组数据标示在以 $C$ 为横坐标、$T^{2}$ 为纵坐标的坐标纸上，即图中用“$\times$”表示的点。
\begin{enumerate}
	\item
	$T$、$L$、$C$ 的关系为\tkanswer[3cm]{$T=2\pi\sqrt{LC}$}。
	\item
	根据图中给出的数据点作出 $T^{2}$ 与 $C$ 的关系图线。
	\item
	求得的 $L$ 值是\tkanswer[3.5cm]{$0.0351\UH\sim0.0389\UH$}。
\end{enumerate}
\onepicture[align=right, w=9cm]{figs/18.png}

%% answer: （1）$T=2\pi\sqrt{LC}$；（2）如图；（3）$0.0351\UH\sim0.0389\UH$
\jdanswer{
\begin{enumerate}
	\item
	$T=2\pi\sqrt{LC}$

	\item
	如图

	\item
	$0.0351\UH\sim0.0389\UH$
\end{enumerate}
}

%% memo:
\memoanswer{
（1）LC 振荡电路的周期公式是 $T=2\pi\sqrt{LC}$。

（2）过点画线时，应使直线通过尽量多的点，不在直线的点均匀地分布在直线两侧，如图所示。

\onepicture[w=9cm]{figs/18_an.png}

（3）从图线上读取一组合适 $T^{2}$、$C$ 数据，由 $L=\frac{T^{2}}{4\pi^{2}C}$ 可算出 $L$ 值，由于所画直线位置存在偏离，故算出 $L$ 值在一定误差范围内就算正确。
}

\item
%% number: 19
%% paperName: 1998年全国高考第 19 题 5 分
%% typeId: 4
%% type: 实验题
%% chapter: 运动和力
%% point: 验证牛顿第二定律
%% method: 无
%% score: 5
%% degree: 650
%% duplicateId: 0
%% body:
某学生做“验证牛顿第二定律”的实验在平衡摩擦力时，把长木板的一端垫得过高，使得倾角偏大。他所得到的 $a-F$ 关系可用下列哪根图线表示？图中 $a$ 是小车的加速度，$F$ 是细线作用于小车的拉力。答：\tkanswer[2cm]{C}。
\fourpicture[labelA=1998全国19a, labelB=1998全国19b, labelC=1998全国19c, labelD=1998全国19d, numA=A, numB=B, numC=C, numD=D, labelprefix={}, h=3.0cm, align=right]
{figs/19a.png}
{figs/19b.png}
{figs/19c.png}
{figs/19d.png}

%% answer: C
\jdanswer{
C
}

%% memo:
\memoanswer{
长木板一端垫得过高，倾角过大时，下滑分力大于摩擦力，没有沙桶的牵引力时，小车所受合外力不为零，便有加速度，从 $a-F$ 图象看，即 $F=0$ 时，$a\neq0$，故 C 正确。
}

\item
%% number: 20
%% paperName: 1998年全国高考第 20 题 6 分
%% typeId: 4
%% type: 实验题
%% chapter: 电路
%% point: 电路实验
%% method: 无
%% score: 6
%% degree: 450
%% duplicateId: 0
%% body:
用图示的电路测定未知电阻 $R_{x}$ 的值。图中电源电动势未知，电源内阻与电流表的内阻均可忽略不计，$R$ 为电阻箱。
\begin{enumerate}
	\item
	若要测得 $R_{x}$ 的值，$R$ 至少需要取\tkanswer[1.5cm]{2}个不同的数值。
	\item
	若电流表每个分度表示的电流值未知，但指针偏转角度与通过的电流成正比，则在用此电路测 $R_{x}$ 时，$R$ 至少需取\tkanswer[1.5cm]{2}个不同的数值。
	\item
	若电源内阻不可忽略，能否应用此电路测量 $R_{x}$？答：\tkanswer[2cm]{不能}。
\end{enumerate}
\onepicture[align=right, w=5cm]{figs/20.png}

%% answer: （1）2；（2）2；（3）不能
\jdanswer{
\begin{enumerate}
	\item
	2

	\item
	2

	\item
	不能
\end{enumerate}
}

%% memo:
\memoanswer{
（1）调整电阻箱使其阻值为 $R_{1}$，读出电流表示数 $I_{1}$，则 $I_{1}=\frac{E}{R_{1}+R_{x}}$，式中 $E$、$R_{x}$ 未知，须再调整电阻箱成另一值 $R_{2}$，读出电流 $I_{2}$，则 $I_{2}=\frac{E}{R_{2}+R_{x}}$，解两式可得 $R_{x}$；

（2）仍需取两个不同 $R$ 值，同（1）的方法，只不过此次读取的电流不是具体数值，而可以用分度来表示；

（3）电源内阻不可忽略时，这样测出来的阻值是 $R_{x}+r$，因电源内阻 $r$ 未知，故不能确定 $R_{x}$ 的值。
}

\item
%% number: 21
%% paperName: 1998年全国高考第 21 题 9 分
%% typeId: 6
%% type: 计算题
%% chapter: 曲线运动
%% point: 平抛运动
%% method: 无
%% score: 9
%% degree: 450
%% duplicateId: 0
%% body:
宇航员站在一星球表面上的某高处，沿水平方向抛出一个小球。经过时间 $t$，小球落到星球表面，测得抛出点与落地点之间的距离为 $L$。若抛出时的初速增大到 2 倍，则抛出点与落地点之间的距离为 $\sqrt{3}L$。已知两落地点在同一水平面上，该星球的半径为 $R$，万有引力常数为 $G$。求该星球的质量 $M$。

%% answer: $M=\frac{2\sqrt{3}LR^{2}}{3Gt^{2}}$
\jdanswer{
$M=\frac{2\sqrt{3}LR^{2}}{3Gt^{2}}$
}

%% memo:
\memoanswer{
设抛出点的高度为 $h$，第一次平抛的水平射程为 $x$，则有 $x^{2}+h^{2}=L^{2}$ ①

由平抛运动规律得知，当初速增大到 2 倍，其水平射程也增大到 $2x$，可得 $(2x)^{2}+h^{2}=(\sqrt{3}L)^{2}$ ②

由①、②解得 $h=\frac{L}{\sqrt{3}}$

设该星球上的重力加速度为 $g$，由平抛运动的规律，得 $h=\frac{1}{2}gt^{2}$

由万有引力定律与牛顿第二定律，得 $\frac{GMm}{R^{2}}=mg$

式中 $m$ 为小球的质量，联立以上各式，解得 $M=\frac{2\sqrt{3}LR^{2}}{3Gt^{2}}$
}

\item
%% number: 22
%% paperName: 1998年全国高考第 22 题 9 分
%% typeId: 6
%% type: 计算题
%% chapter: 气体性质
%% point: 气体连接体
%% method: 无
%% score: 9
%% degree: 450
%% duplicateId: 0
%% body:
活塞把密闭气缸分成左、右两个气室，每室各与 U 形管压强计的一臂相连。压强计的两臂截面处处相同。U 形管内盛有密度为 $\rho=7.5\times10^{2}\Ukgmc$ 的液体。开始时左、右两气室的体积都为 $V_{0}=1.2\times10^{-2}\Umc$，气压都为 $p_{0}=4.0\times10^{3}\UPa$，且液体的液面处在同一高度，如图所示。现缓缓向左推进活塞，直到液体在 U 形管中的高度差 $h=40\Ucm$。求此时左、右气室的体积 $V_{1}$、$V_{2}$。假定两气室的温度保持不变。计算时可以不计 U 形管和连接管道中气体的体积。
\onepicture[align=right, w=5cm]{figs/22.png}

%% answer: $V_{1}=8.0\times10^{-3}\Umc$，$V_{2}=1.6\times10^{-2}\Umc$
\jdanswer{
$V_{1}=8.0\times10^{-3}\Umc$，$V_{2}=1.6\times10^{-2}\Umc$
}

%% memo:
\memoanswer{
以 $p_{1}$、$V_{1}$ 表示压缩后左室气体的压强和体积，$p_{2}$、$V_{2}$ 表示这时右室气体的压强和体积；$p_{0}$、$V_{0}$ 表示初态两室气体的压强和体积，缓慢推进活塞时，气体做等温变化，由玻意耳定律得

对左室气体有 $p_{1}V_{1}=p_{0}V_{0}$ ①

对右室气体有 $p_{2}V_{2}=p_{0}V_{0}$ ②

由体积不变得 $V_{1}+V_{2}=2V_{0}$ ③

由压强平衡得 $p_{1}-p_{2}=\Delta p=\rho gh$ ④

解以上四式得：$V_{1}^{2}-\frac{2(p_{0}+\Delta p)V_{0}}{\Delta p}V_{1}+\frac{2p_{0}V_{0}^{2}}{\Delta p}=0$ ⑤

解方程并选择物理意义正确的解得到 $V_{1}=\frac{V_{0}}{\Delta p}\left(p_{0}+\Delta p-\sqrt{p_{0}^{2}+\Delta p^{2}}\right)$

代入数值，得：$V_{1}=8.0\times10^{-3}\Umc$ ⑥

$V_{2}=2V_{0}-V_{1}=1.6\times10^{-2}\Umc$ ⑦
}

\item
%% number: 23
%% paperName: 1998年全国高考第 23 题 11 分
%% typeId: 6
%% type: 计算题
%% chapter: 磁场
%% point: 洛伦兹力
%% method: 无
%% score: 11
%% degree: 450
%% duplicateId: 0
%% body:
如图所示，在 $x$ 轴上方有垂直于 $xy$ 平面向里的匀强磁场，磁感应强度为 $B$；在 $x$ 轴下方有沿 $y$ 轴负方向的匀强电场，场强为 $E$。一质量为 $m$，电量为 $-q$ 的粒子从坐标原点 $O$ 沿着 $y$ 轴正方向射出。射出之后，第三次到达 $x$ 轴时，它与点 $O$ 的距离为 $L$。求此粒子射出时的速度 $v$ 和运动的总路程 $s$（重力不计）。
\onepicture[align=right, w=5cm]{figs/23.png}

%% answer: $v=\frac{qBL}{4m}$；$s=\frac{\pi L}{2}+\frac{qB^{2}L^{2}}{16mE}$
\jdanswer{
\begin{enumerate}
	\item
	$v=\frac{qBL}{4m}$

	\item
	$s=\frac{\pi L}{2}+\frac{qB^{2}L^{2}}{16mE}$
\end{enumerate}
}

%% memo:
\memoanswer{
粒子运动路线如图示有 $L=4R$ ①

粒子初速度为 $v$，由洛伦兹力提供向心力得 $qvB=m\frac{v^{2}}{R}$ ②

由①、②式可算得 $v=\frac{qBL}{4m}$ ③

\onepicture[w=7cm]{figs/23_an.png}

设粒子进入电场作减速运动的最大路程为 $l$，加速度为 $a$，由牛顿第二定律和运动学公式得 $qE=ma$ ④

$v^{2}=2al$ ⑤

粒子运动的总路程 $s=2\pi R+2l$ ⑥

由①、②、④、⑤、⑥式，得 $s=\frac{\pi L}{2}+\frac{qB^{2}L^{2}}{16mE}$ ⑦
}

\item
%% number: 24
%% paperName: 1998年全国高考第 24 题 12 分
%% typeId: 6
%% type: 计算题
%% chapter: 光学
%% point: 透镜
%% method: 无
%% score: 12
%% degree: 450
%% duplicateId: 0
%% body:
如图所示，$L$ 为薄凸透镜，点光源 $S$ 位于 $L$ 的主光轴上，它到 $L$ 的距离为 $36\Ucm$；$M$ 为一与主光轴垂直的挡光圆板，其圆心在主光轴上，它到 $L$ 的距离为 $12\Ucm$；$P$ 为光屏，到 $L$ 的距离为 $30\Ucm$。现看到 $P$ 上有一与挡光板同样大小的圆形暗区 $ab$。求透镜的焦距。
\onepicture[align=right, w=9cm]{figs/24.png}

%% answer: $f_{2}=12\Ucm$
\jdanswer{
$f_{2}=12\Ucm$
}

%% memo:
\memoanswer{
光屏上的暗区是由于挡光圆板挡住部分光线而形成的。因而从点光源 $S$ 经挡光圆板边缘（譬如图中的 $c$ 点）射到透镜上 $H$ 点的光线 $ScH$，经透镜折射后，出射光线应经过暗区的边缘的某点。这点可以是暗区的边缘点 $a$，也可以是暗区的另一边缘点 $b$。也就是说符合要求的像点有两点：$S_{1}'$、$S_{2}'$。

\onepicture[w=10cm]{figs/24_an.png}

先求与像点 $S_{1}'$ 相应的焦距 $f_{1}$。设 $r$ 表示圆板与暗区的半径，以 $u$ 表示物距，$v_{1}$ 表示像距，则有 $\frac{\overline{OH}}{r}=\frac{u}{u-l_{1}}$ ①

$\frac{\overline{OH}}{r}=\frac{v_{1}}{v_{1}-l_{2}}$ ②

由成像公式，得 $\frac{1}{u}+\frac{1}{v_{1}}=\frac{1}{f_{1}}$ ③

解①、②、③式得 $f_{1}=25.7\Ucm$ ④

再求与像点 $S_{2}'$ 相应的焦距 $f_{2}$，以 $v_{2}$ 表示像距，$\frac{\overline{OH}}{r}=\frac{v_{2}}{l_{2}-v_{2}}$ ⑤

由成像公式，得 $\frac{1}{u}+\frac{1}{v_{2}}=\frac{1}{f_{2}}$ ⑥

解①、⑤、⑥式得 $f_{2}=12\Ucm$ ⑦
}

\item
%% number: 25
%% paperName: 1998年全国高考第 25 题 12 分
%% typeId: 6
%% type: 计算题
%% chapter: 运动和力
%% point: 动量和能量
%% method: 无
%% score: 12
%% degree: 250
%% duplicateId: 0
%% body:
一段凹槽 A 倒扣在水平长木板 C 上，槽内有一小物块 B，它到槽两内侧的距离均为 $\frac{1}{2}L$，如图所示。木板位于光滑水平的桌面上，槽与木板间的摩擦不计，小物块与木板间的摩擦系数为 $\mu$。A、B、C 三者质量相等，原来都静止。现使槽 A 以大小为 $v_{0}$ 的初速向右运动，已知 $v_{0}<\sqrt{2\mu gL}$。当 A 和 B 发生碰撞时，两者速度互换。求：
\begin{enumerate}
	\item
	从 A、B 发生第一次碰撞到第二次碰撞的时间内，木板 C 运动的路程。
	\item
	在 A、B 刚要发生第四次碰撞时，A、B、C 三者速度的大小。
\end{enumerate}
\onepicture[align=right, w=8cm]{figs/25.png}

%% answer: （1）$s=l-\frac{v_{0}^{2}}{4\mu g}$；（2）$v_{A}=\frac{1}{4}v_{0}$，$v_{B}=v_{C}=\frac{3}{8}v_{0}$
\jdanswer{
\begin{enumerate}
	\item
	$s=l-\frac{v_{0}^{2}}{4\mu g}$

	\item
	$v_{A}=\frac{1}{4}v_{0}$，$v_{B}=v_{C}=\frac{3}{8}v_{0}$
\end{enumerate}
}

%% memo:
\memoanswer{
（1）由题意知，A 和 B 碰撞时，两者速度互换。A 与 B 刚发生第一次碰撞后，A 停下不动，B 以初速 $v_{0}$ 向右运动。由于摩擦，向右作匀减速运动，而 C 向右作匀加速运动，两者速率逐渐接近。设 B、C 达到相同速度 $v_{1}$ 时 B 移动的路程为 $s_{1}$。设 A、B、C 质量皆为 $m$，由动量守恒定律，得 $mv_{0}=2mv_{1}$ ①

由功能关系，得 $\mu mgs_{1}=\frac{1}{2}mv_{0}^{2}-\frac{1}{2}mv_{1}^{2}$ ②

由①得 $v_{1}=\frac{1}{2}v_{0}$

代入②式，得 $s_{1}=\frac{3v_{0}^{2}}{8\mu g}$

根据条件 $v_{0}<\sqrt{2\mu gl}$，得 $s_{1}<\frac{3}{4}l$ ③

可见，在 B、C 达到相同速度 $v_{1}$ 时，B 尚未与 A 发生第二次碰撞，B 与 C 一起将以 $v_{1}$ 向右匀速运动一段距离 $(l-s_{1})$ 后才与 A 发生第二次碰撞。

设 C 的速度从零变到 $v_{1}$ 的过程中，C 的路程为 $s_{2}$。由功能关系，得 $\mu mgs_{2}=\frac{1}{2}mv_{1}^{2}$ ④

解得 $s_{2}=\frac{v_{0}^{2}}{8\mu g}$

因此在第一次到第二次碰撞间 C 的路程为 $s=s_{2}+l-s_{1}=l-\frac{v_{0}^{2}}{4\mu g}$ ⑤

（2）由上面讨论可知，在刚要发生第二次碰撞时，A 静止，B、C 的速度均为 $v_{1}$。刚碰撞后，B 静止，A、C 的速度均为 $v_{1}$。由于摩擦，B 将加速，C 将减速，直至达到相同速度 $v_{2}$。由动量守恒定律，得 $mv_{1}=2mv_{2}$ ⑥

解得 $v_{2}=\frac{1}{2}v_{1}=\frac{1}{4}v_{0}$

因 A 的速度 $v_{1}$ 大于 B 的速度 $v_{2}$，故第三次碰撞发生在 A 的左壁。刚碰撞后，A 的速度变为 $v_{2}$，B 的速度变为 $v_{1}$，C 的速度仍为 $v_{2}$。由于摩擦，B 减速，C 加速，直至达到相同速度 $v_{3}$。由动量守恒定律，得 $mv_{1}+mv_{2}=2mv_{3}$ ⑦

解得 $v_{3}=\frac{3}{8}v_{0}$

故刚要发生第四次碰撞时，A、B、C 的速度分别为 $v_{A}=v_{2}=\frac{1}{4}v_{0}$ ⑧

$v_{B}=v_{C}=v_{3}=\frac{3}{8}v_{0}$ ⑨
}

\end{enumerate}

\end{document}
